\documentclass[
10pt,
a4paper,
oneside,
headinclude,footinclude,
BCOR5mm,
]{article}

\input{structure.tex} % Include the structure.tex file which specified the document structure and layout

\hyphenation{Fortran hy-phen-ation} % Specify custom hyphenation points in words with dashes where you would like hyphenation to occur, or alternatively, don't put any dashes in a word to stop hyphenation altogether

\begin{document}

\begin{titlepage}
  \centering
  {\scshape\LARGE Universidad Nacional Autónoma de México \par}
      \vspace{1cm}
	     {\scshape\Large Facultad de Ciencias\par}
	     \vspace{1.2cm}
             \begin{center}
	       \includegraphics[scale=.2]{../imagenes/UNAM.eps}
	     \end{center}
	     \vspace{.5 cm}

	     {\huge\bfseries Proyecto 1: \par}
	     {\huge\bfseries Travelling Salesman Problem \par}
	     \vspace{0.5cm}

	     {\Large\itshape Emiliano Galeana Araujo\\314032324\par}
	     \vfill
	     \vspace{0.5cm}

             Seminario de Ciencias de la computación B\par
             Heuríticas de Optimización Combinatoria\par\par
	     \textsc{Dr. Canek Peláez Valdés}\\
             \textsc{L. en C.C. Víctor Zamora Gutiérrez}\\
	     \vspace{0.1cm}
	            {\large \today \par}
\end{titlepage}


%----------------------------------------------------------------------------------------
%	HEADERS
%----------------------------------------------------------------------------------------

\renewcommand{\sectionmark}[1]{\markright{\spacedlowsmallcaps{#1}}} % The header for all pages (oneside) or for even pages (twoside)
%\renewcommand{\subsectionmark}[1]{\markright{\thesubsection~#1}} % Uncomment when using the twoside option - this modifies the header on odd pages
\lehead{\mbox{\llap{\small\thepage\kern1em\color{halfgray} \vline}\color{halfgray}\hspace{0.5em}\rightmark\hfil}} % The header style

\pagestyle{scrheadings} % Enable the headers specified in this block

%----------------------------------------------------------------------------------------
%	TABLE OF CONTENTS & LISTS OF FIGURES AND TABLES
%----------------------------------------------------------------------------------------

\setcounter{tocdepth}{2} % Set the depth of the table of contents to show sections and subsections only

\tableofcontents

\listoffigures

\section*{Abstract}

Se conocen muchos problemas de optimización. El objetivo de este proyecto es el
\TSP (TSP) pero con la particularidad de que no se busca un ciclo, sino un camino
que pase por \textit{n} ciudades sin repetir ninguna.\\
Dada una base de datos con ciudades, distancias entre ellas y conexiones entre
ciudades. ¿Cuál es el mínimo costo que podemos encontrar para ejemplares de 40 y
150 ciudades? Se emplea la heurística de recocido simulado, la cual sirve para
resolver problemas de optimización. Los resultados arrojadas por la heurística se
cree son los óptimos, sin embargo no se pueden demostrar.

\newpage


\section{Introduction}

A statement requiring citation \cite{Figueredo:2009dg}.

\section{Lenguaje}

Para la implementación de la heurística se eligió como lenguaje
\texttt{Go(Golang)} el cual \textit{es un lenguaje open source que hace más
  sencillo construir software de manera simple, confiable y eficiente.
  \cite{Golang01}} Lo elegí básicamente porque había escuchado que era bueno y
no era muy complicado de usar y porque nunca había hecho algo en \GO. La
curva de aprendizaje fue complicada, pues empecé con algunas funciones en una
carpeta hecha para el curso, en donde tenía al \textit{main} ahí mismo. Al
momento de pasar al main a otro archivo comenzaron los problemas y después de
buscar en internet descubrí que tenía que migrar mis archivos a la carpeta que
\GO crea en el \textit{home}.\\
Honestamente creí que lo más complicado sería poder conectar mi programa a la
base de datos, pero fue muy fácil, hay muchos tutoriales en internet sobre
\GO que me ayudaron con eso y con otras cosas del lenguaje como las firmas de las
funciones y lo mejor de todo, que me sirvió para algunas funciones de este
proyecto. Que puedes tener funciones que regresan más de un tipo.\\
Los contras es que no podía copiar un array tan sencillo, que aún después de
crear un array, podía seguir insertando, que fue complicada la diferencia entre
\texttt{a = b} y \texttt{a := b}. Y obviamente el hecho de tener que estar en la
carpeta generada en el \textit{home}.
Los pros fueron los tipos de regreso, que cuando importabas una biblioteca si no
la usabas no te dejaba compilar el programa, igual que cuando creabas una
variable y esta no era usada no te dejaba compilar, y que una vez que más o menos
le agarras la onda ya no es tan complicado hacer cosas sencillas en el lenguaje.

\section{Packages}

\begin{enumerate}[noitemsep]
\item main
\item funciones
\item argumentos
\end{enumerate}

\subsection{Main}

En este ``package'' lo que se implementó es la función main, y se importan los
``package'' \textit{``fmt'', ``os'', ``math/rand'', ``funciones'', ``argumentos''
}.

\paragraph{Función main} Es la función principal, se encarga de verificar la
entrada esto quiere decir, que se le pasen los argumentos correspondientes.
También se crea un ``struct'' Ciudades y Solucion que son llamados de
``funciones''.

\subsection{Argumentos}
En este ``package'' se encuentra el archivo ``leer\_ciudades.go''.
\paragraph{Función Leer} Es la función que recibe los argumentos pasados por el
main. Una seed y un archivo con las ciudades que se buscan para el camino, este
archivo tiene que ser de una línea y con los índices de las ciudades separados
por coma (,). Los índices de las ciudades se pueden encontrar en la base de
datos. La función regresa un \texttt{[]int} con los índices
de las ciudades, regresa la semilla que es el segundo argumento en tipo
\texttt{int}, y el nombre del archivo que se leyó, esto con la finalidad de
escribir un archivo con los datos obtenidos y se tenga conocimiento de a qué
resultados hacen referencia.

\subsection{Funciones}
En este ``package'' se encuentran los archivos ``funciones.go'',
``operaciones.go'' y ``sql.go'', así como sus respectivos tests.

\subsubsection{funciones.go}

\paragraph{Función PrintCiudad}
\paragraph{Función PrintData}
\paragraph{Función totalAristas}
\paragraph{Función GetNormalizador}
\paragraph{Función FunCosto}
\paragraph{Función FunCostoSolucion}
\paragraph{Función porcentajeAceptados}
\paragraph{Función busquedaBinaria}
\paragraph{Función TemperaturaInicial}
\paragraph{Función calculaLote}
\paragraph{Función AceptacionPorUmbrales}
\paragraph{Función printSol}
\paragraph{Función NewSolucion}
\paragraph{Función NewCiudades}


\subsubsection{operaciones.go}
En este archivo se implementan todas las funciones que no tienen que ver tanto
con la heurística o que como su nombre lo indica, son operaciones básicas.

\paragraph{Función radianes} Cambia una coordenada (\texttt{float64}) en su
representación en radianes (\texttt{float64}).
\paragraph{Función distanciaNatural} Calcula la distancia en la vida real entre
dos índices de, ciudades, no es la distancia exacta. Hace uso del archivo
``sql.go''.
\paragraph{Función obtenerA} Regresa el resultado de la siguiente fórmula:
\begin{equation}
  A = \sin(\frac{lat(v)- lat(u)}{2})^{2} + \cos(lat(u)) \times \cos(lat(v)) \times
  \sin(\frac{lon(v) - lon(u)}{2})^{2}
\label{eq:refname2}
\end{equation}

\paragraph{Función pesoAumentado} Dados dos índices de ciudades y la máxima distancia,
regresa la distancia entre las dos ciduades multiplicada por la máxima distancia.
\paragraph{Función CopiarCiudades} Copia un \texttt{[]int} y regresa la copia.
\paragraph{Función swap} Dados dos índices, y un \texttt{[]int}, cambia los índices en
una copia del \texttt{[]int}.
\paragraph{Función vecino} Dado un \texttt{[]int}, que representa el acomodo de los
índices de las ciudades, regresa un vecino de esta representación, eso es, el
intercambio de dos índices aleatorios.


\subsubsection{sql.go}
En este archivo se hace la conexión con la base de datos, así como los queries
necesarios durante la implementación de la heurística.

\paragraph{obtenerLatLon} Es la función encargada de regresar una tupla de
\texttt{float64} la tupla representa a la latitud y longitud de las ciudades en
el mundo.

\paragraph{completa} Esta función recibe un \texttt{[]int} el cual contiene a los
índices de las ciudades y regresa un \texttt{[][]float64} el cual representa a
todas las conexiones que se quieren para una instancia de nuestro problema.

\section{Resultados}


\bibliography{mybib}{}
\bibliographystyle{plain}

\end{document}
